\chapter*{Conclusion générale et perspectives}
\addcontentsline{toc}{chapter}{Conclusion générale et perspectives}
\markboth{Conclusion générale et perspectives}{}

Ce projet avait pour objectif de concevoir une Modern Data Platform capable d'automatiser l'ingestion,
la transformation et la modélisation des données, et d'offrir aux décideurs une interface leur
permettant d'interroger ces données en langage naturel.

Cet objectif est atteint. Les données transactionnelles de Wide World Importers sont extraites de
SQL Server par \textbf{DLT} de manière incrémentale et chargées dans \textbf{Snowflake}, où elles sont
organisées selon une \textbf{architecture Médaillon}. La couche Bronze conserve une copie fidèle et
réconciliée de la source ; la couche Silver applique des règles de nettoyage homogènes et factorisées ;
la couche Gold expose un schéma en étoile complété de dimensions historisées en SCD de type 2. Les
transformations, développées avec \textbf{dbt}, sont versionnées, documentées et couvertes par plus de
140 tests. L'ensemble est orchestré par \textbf{Dagster}, qui offre un lignage de bout en bout, une
exécution quotidienne avec reprise automatique, des contrôles de fraîcheur et des alertes.

L'application \textbf{Streamlit}, adossée à \textbf{Snowflake Cortex Analyst} et à un modèle sémantique
riche, répond correctement aux questions de l'énoncé et à des questions plus élaborées, y compris des
questions de suivi, en restituant le SQL généré, un graphique adapté et le tableau des résultats. Sa
conception place la sécurité au premier plan : contrôle de lecture seule, rôle aux droits minimaux et
absence d'identifiants dans le code. Déployée aussi dans \textbf{Streamlit in Snowflake} et validée par
une chaîne d'\textbf{intégration continue} sur clone zéro copie, la plateforme dépasse le stade du
prototype pour s'approcher d'une solution industrialisée.

Sur le plan personnel, ce projet m'a permis de mettre en pratique l'ensemble de la chaîne de valeur de
la donnée --- ingénierie, modélisation dimensionnelle, qualité, orchestration, sécurité, DevOps --- et
de mesurer l'apport, mais aussi les exigences, de l'intelligence artificielle générative appliquée à
l'analyse de données : un LLM n'est fiable que s'il s'appuie sur un modèle de données propre et sur
une sémantique métier explicitement décrite.

\section*{Perspectives}

Plusieurs prolongements sont envisageables :
\begin{itemize}
  \item \textbf{Enrichissement du modèle sémantique} : nouvelles requêtes vérifiées issues des questions
        réelles des utilisateurs, publication sous forme de \emph{semantic view} Snowflake, et
        couverture des historisations SCD2 et des températures d'entrepôt ;
  \item \textbf{Retour utilisateur} : collecte des appréciations (réponse correcte ou non) pour mesurer
        en continu la précision du Text-to-SQL ;
  \item \textbf{Ingestion en quasi temps réel} : exploitation de la capture des changements (CDC) de
        SQL Server ou d'un ordonnancement plus fréquent ;
  \item \textbf{Gouvernance} : masquage dynamique des données personnelles, politiques d'accès par ligne
        (\emph{row access policies}) selon le territoire commercial de l'utilisateur, catalogue de
        données ;
  \item \textbf{Observabilité} : tableau de bord de suivi des coûts Snowflake, des durées d'exécution et
        des résultats des tests ;
  \item \textbf{Restitution} : tableaux de bord prédéfinis (Power BI, Streamlit) pour les indicateurs
        récurrents, l'interface conversationnelle restant dédiée aux questions ad hoc.
\end{itemize}
